%!TEX root = ../thesis.tex
\chapter*{Abstract}

There has been an increase in the deployment of underwater sensor networks for environmental monitoring. This trend has produced a need for database systems capable of managing large volumes of time-series data efficiently. The SFI Smart Ocean project looks to create an autonomous marine observation platform which integrates underwater sensor networks with cloud-based infrastructure for use in long-term environment monitoring and industrial marine operations. As these deployments expand, choosing an appropriate database technology becomes a vital architectural decision because of the differing performance characteristics of the available database systems. This study examined the suitability of TimescaleDB, MongoDB and InfluxDB for the management of large-scale time-series data generated by representative Smart Ocean workloads. The research designed and implemented a benchmarking framework that is also reproducible. The framework created a controlled environment to evaluate selected database technologies under identical experimental conditions. Synthetic observation datasets representative of Smart Ocean sensor data was generated in progressively increasing sizes and they were loaded into each database system. Storage efficiency, ingestion performance and query performance were studied. Ingestion performance was assessed using observation throughput, execution time and data throughput. Query performance was evaluated using latest-value query latency and storage efficiency was measured via database storage growth and average storage consumption for each observation. The benchmark results showed major performance differences among the three database systems. MongoDB had the shortest ingestion execution times, the highest data throughput and observation throughput and the lowest latest-value query latency and storage footprint. InfluxDB demonstrated more balanced performance across the evaluated metrics, maintaining stable query and ingestion behaviors with moderate requirements for storage. TimescaleDB successfully supported all benchmark workloads, but generally showed greater query latency, higher execution times and larger storage requirements than the other evaluated database systems. The study concludes that the architecture of each database has a major influence on the performance characteristics of large-scale time-series data management systems. All three databases demonstrated different performance trade-offs associated with their architectures. The benchmarking framework and empirical findings are contributing evidence supporting database selection decisions for Smart Ocean and similar environmental monitoring applications. 


\textbf{Keywords:} Time-series data, database benchmarking, Smart Ocean, underwater sensor networks, MongoDB, InfluxDB, TimescaleDB, ingestion performance, query latency, storage efficiency.




